% Build from the repository root, twice, after creating the output directory:
% pdflatex -interaction=nonstopmode -halt-on-error -output-directory=verification/anisotropic-geometry-20261026/latex tex/result/anisotropic_full_moments.tex
\documentclass[11pt,a4paper]{article}
\usepackage[T2A]{fontenc}
\usepackage[utf8]{inputenc}
\usepackage[english,russian]{babel}
\usepackage[margin=21mm]{geometry}
\usepackage{amsmath,amssymb,amsthm,booktabs,microtype}
\usepackage[hidelinks]{hyperref}
\usepackage{xurl}
\setlength{\parindent}{0pt}
\setlength{\parskip}{4pt}
\setlength{\emergencystretch}{3em}
\newcommand{\R}{\mathbb R}
\newcommand{\tr}{\operatorname{tr}}
\newcommand{\spn}{\operatorname{span}}
\newcommand{\sym}{\operatorname{sym}}
\newcommand{\rank}{\operatorname{rank}}
\newcommand{\diag}{\operatorname{diag}}
\newcommand{\Hess}{\operatorname{Hess}}
\newcommand{\grad}{\operatorname{grad}}
\newcommand{\EDR}{\operatorname{EDR}}
\newcommand{\norm}[1]{\lVert#1\rVert}
\newcommand{\inner}[2]{\langle#1,#2\rangle_F}
\newcommand{\code}[1]{\nolinkurl{#1}}
\newtheorem{theorem}{Теорема}[section]
\newtheorem{proposition}[theorem]{Утверждение}
\newtheorem{corollary}[theorem]{Следствие}
\theoremstyle{remark}
\newtheorem{remark}[theorem]{Замечание}
\title{Анизотропная геометрия и ridge\newline для полных локальных EDR-моментов}
\author{Математическая переработка черновика}
\date{26 октября 2026 года}
\begin{document}
\maketitle

\textbf{Основной вывод.}
Для исследования оптимизации разумно оставить евклидово ортонормированный
базис и исходный статистический критерий, изменяя только метрику касательного
шага. Полная форма локальной кривизны учитывает одновременно ковариации
признаков и силу локальных индексов. Её совпадение с половиной гессиана
является точным результатом при нулевом остатке, отсутствии ridge и
положительной определённости локальных систем. Для действующего ridge
нужны другие формулы; улучшение EDR из изменения метрики не следует.

В этой записке исправлены ограничения по рангу, выведены точный ridge-гессиан
и две разные формы Гаусса--Ньютона, упрощено применение двусторонней метрики
и получена условная локальная оценка обычной EDR. Все выводы относятся к
\emph{одному внутреннему решению с фиксированными окнами}. Новых полных
ADP-fit и сравнения качества восстановления здесь нет.

\section{Фиксируем задачу и связь с текущим кодом}

Пусть $1\le m<d$, $B\in\R^{d\times m}$, $B^\top B=I_m$,
$\pi_j>0$, $\sum_j\pi_j=1$. Внутри одного решения фиксированы
симметричные $C_j\succeq0$, $c_j\in\operatorname{range}(C_j)$ и
$\tau\ge0$. Во всей записке $C_j,c_j$ обозначают уже \emph{спектрально
усечённые} $\widehat C_j,\widehat c_j$ из \code{full_moments_ridge.tex}.
Положим
\begin{equation}\label{eq:profile}
 U_j=C_j^{1/2},\qquad y_j=C_j^{\dagger/2}c_j,\qquad
 F_\tau(B)=\sum_j\pi_j\min_g
 \left\{\norm{y_j-U_jBg}^2+\tau\norm{g}^2\right\}.
\end{equation}
При $\tau>0$, а также при $\tau=0$ и $B^\top C_jB\succ0$,
\begin{equation}\label{eq:normal}
 A_j=B^\top C_jB+\tau I_m,\qquad
 g_j=A_j^{-1}B^\top c_j,\qquad r_j=C_jBg_j-c_j.
\end{equation}
Используемый далее $r_j$ --- \emph{моментный} остаток; он не равен
отбелённому остатку $y_j-U_jBg_j$. Градиент обычного продолжения
\eqref{eq:profile} во внешнее пространство равен
\begin{equation}\label{eq:gradient}
 G_E=2\sum_j\pi_jr_jg_j^\top,\qquad
 P=I_d-BB^\top,\qquad G_H=PG_E.
\end{equation}

Проверено по \code{experiments/coupled_penalties.py},
\code{experiments/full_ridge_fast.py},
\code{experiments/full_ridge_batched.py}: экспериментальный
\code{full} и его пакетный вариант используют этот профиль, спектральный
порог $10^{-12}$ и
\[
 \tau=0.01\,\operatorname{median}_j
 \norm{C_j^{1/2}B_{\rm in}}_F^2/m.
\]
Этот скаляр фиксируется до оптимизации. Пакетная реализация использует
QR-карту с фиксированным начальным дополнением и обычный
\code{scipy.optimize.minimize(method="L-BFGS-B")}; это не реализация
переменной римановой метрики. Она возвращает строковый базис $B^\top$.
Защитная ветка при очень малом ridge возвращается к отбелённому профилю.
Ни эта ветка, ни штатный \code{ADP/solver/grassman.py} здесь не меняются.

\textbf{Исправление утверждения о ранге.}
Для локальных весов $a_{ji}>0$ на $n_{+,j}$ точках центрированная ковариация
имеет ранг не выше $\min(d,n_{+,j}-1)$. Параметр \code{N_loc} задаёт массу
локального ядра, поэтому \code{N_loc=20} не означает $n_{+,j}=20$ и не
даёт границы $\rank C_j\le19$. В существующем артефакте
\code{verification/full-ridge-batched-20261025/native/m2-974100-direct.json},
\code{fits[0].stages[0].penalty.full_moments}, записана гистограмма
с рангами от $0$ до $25$ и $209$ полноранговыми окнами из $300$.
Это проверка сохранённого результата, а не новый запуск.
Следовательно, применимость теорем проверяется по фактическим $C_j$ и
$A_j$, а не по \code{N_loc}.

\section{Отдельно: нормировка базиса, штраф и метрика}

Любое $m$-мерное подпространство имеет как евклидов ортонормированный базис
$B$, так и $M$-ортонормированный базис
\begin{equation}\label{eq:same-subspace}
 Z=B(B^\top MB)^{-1/2},\qquad Z^\top MZ=I_m,
 \qquad M\succ0.
\end{equation}
При фиксированном $M$ переход между ограничениями допускает те же
подпространства. Преобразование $Q=M^{1/2}Z$ даёт обычный Штифель,
но $\spn(Q)=M^{1/2}\spn(Z)$ обычно является \emph{другим}
подпространством. Поэтому при вычислениях в $Q$ надо преобразовать и
функционал: $F(Z)=F(M^{-1/2}Q)$.

При $\tau=0$ минимум по $g$ инвариантен при $B\mapsto BL$ для любого
обратимого $L$: достаточно заменить $g\mapsto L^{-1}g$.
Положительная определённость $B^\top C_jB$ нужна для гладких производных,
но не для этого равенства значений минимума. При $\tau>0$
штраф $\norm{g}^2$ инвариантен лишь относительно ортогональных $L$.

Для физического вектора локального градиента $v=Zg$ имеем
\[
 Z^\top MZ=I_m\quad\Longrightarrow\quad
 \tau\norm{g}^2=\tau v^\top Mv.
\]
Таким образом, неизменённый коэффициентный ridge после изменения
ограничения меняет статистическую задачу. Чтобы сохранить прежний
евклидов штраф, на обобщённом Штифеле надо использовать
\begin{equation}\label{eq:invariant-ridge}
 \tau g^\top Z^\top Zg=\tau\norm{Zg}^2.
\end{equation}
Тогда локальная система и внешний градиент имеют вид
\begin{align}\label{eq:invariant-gradient}
 g_j&=[Z^\top(C_j+\tau I_d)Z]^{-1}Z^\top c_j,\nonumber\\
 G_E^{\rm inv}&=2\sum_j\pi_j[(C_j+\tau I_d)Zg_j-c_j]g_j^\top.
\end{align}
Второй член $2\tau\sum_j\pi_jZg_jg_j^\top$ нельзя просто забыть.
На обычном Grassmann его горизонтальная проекция равна нулю; на
обобщённом Штифеле надо дифференцировать именно \eqref{eq:invariant-ridge}.
Для чистого сравнения геометрии величину $\tau$ также следует вычислять
один раз из общего \emph{евклидова} $B_{\rm in}$ до перенормировки.

\section{Обычная EDR остаётся мерой восстановления}

Для произвольного полнорангового $Z$ и истинного ортонормированного
$B_*$ определим
\begin{equation}\label{eq:edr}
 P_E(Z)=Z(Z^\top Z)^{-1}Z^\top,\qquad
 \EDR(Z,B_*)=\frac{\norm{P_E(Z)-B_*B_*^\top}_F}{\sqrt{2m}}
 =\left(\frac1m\sum_{i=1}^m\sin^2\theta_i\right)^{1/2}.
\end{equation}
Именно такая нормировка используется в
\code{experiments/r_sdr/design.py:metrics}.
Для вычислений предпочтительно сначала получить ортонормированный базис
через QR или SVD, не строя явную обратную матрицу.
Проектор $P_M(Z)=Z(Z^\top MZ)^{-1}Z^\top M$ не является евклидовым
симметричным проектором. Его разность нельзя подставлять в
\eqref{eq:edr} под видом прежней EDR.

Например, при $M=\diag(100,1)$, $B_*=e_1$,
$\spn(Z)=\spn(e_1+e_2)$ имеем
$\EDR=1/\sqrt2$, тогда как синус угла после преобразования
$x\mapsto M^{1/2}x$ равен $1/\sqrt{101}$.
Это разные расстояния для одного и того же оценённого подпространства.

\section{Что верно для фиксированного обобщённого Штифеля}

При фиксированном $M\succ0$ и метрике
$\langle U,V\rangle_M=\tr(U^\top MV)$ на $Z^\top MZ=I_m$:
\begin{align}\label{eq:generalized}
 T_Z\operatorname{St}_M&=\{\Delta:Z^\top M\Delta+\Delta^\top MZ=0\},\nonumber\\
 \grad_MF(Z)&=M^{-1}G_E-Z\sym(Z^\top G_E),\nonumber\\
 R_Z(\Delta)&=(Z+\Delta)
 [(Z+\Delta)^\top M(Z+\Delta)]^{-1/2}.
\end{align}
Для корректного касательного $\Delta$ внутренняя матрица равна
$I_m+\Delta^\top M\Delta\succ0$. Формулы черновика здесь верны.
Если критерий инвариантен к правым ортогональным поворотам, задача
спускается на фактор по $O(m)$, то есть на пространство подпространств.
Конструкция согласуется с общим аппаратом \cite{shustin}.

Изменение ограничения и метрики не являются одним и тем же действием.
В частности, метрика, полученная из \eqref{eq:generalized} после
перехода к обычному представителю $B$ того же подпространства, равна
\begin{equation}\label{eq:pullback}
 \langle\Delta_1,\Delta_2\rangle_{\rm pull}
 =\tr\left[\Delta_1^\top S_M(B)\Delta_2(B^\top MB)^{-1}\right],
 \qquad B^\top\Delta_i=0,
\end{equation}
где
\begin{equation}\label{eq:schur}
 S_M(B)=M-MB(B^\top MB)^{-1}B^\top M.
\end{equation}
Действительно, $M$-горизонтальный подъём $\Delta$ в точку
$Z=B(B^\top MB)^{-1/2}$ равен
\[
 [I_d-B(B^\top MB)^{-1}B^\top M]\Delta(B^\top MB)^{-1/2}.
\]
Его $M$-скалярное произведение даёт \eqref{eq:pullback}.
Следовательно, обобщённый Штифель с естественной метрикой и простая
левосторонняя метрика $S_M(B)$ на обычном Grassmann в общем случае
\emph{не совпадают}: у первой есть дополнительный правый множитель.

\section{Точная кривизна идеального критерия}

Предположим
\begin{equation}\label{eq:ideal}
 c_j=C_jB_*a_j,\qquad B_*^\top B_*=I_m,\qquad
 B_*^\top C_jB_*\succ0,\qquad \tau=0.
\end{equation}
Обозначим
\[
 S_j=C_j-C_jB_*(B_*^\top C_jB_*)^{-1}B_*^\top C_j.
\]
\begin{theorem}\label{th:ideal}
Для горизонтального $\Delta$, $B_*^\top\Delta=0$, справедливо
\begin{equation}\label{eq:ideal-hessian}
 \Hess F_0(B_*)[\Delta,\Delta]
 =2\sum_j\pi_j(\Delta a_j)^\top S_j(\Delta a_j).
\end{equation}
Здесь гессиан в критической точке рассматривается как билинейная форма;
она не зависит от выбора римановой метрики.
\end{theorem}
\begin{proof}
Для кривой $B(t)=B_*+t\Delta+O(t^2)$ и
$g_j(t)=a_j+tu_j+O(t^2)$ остаток первого порядка равен
$-tC_j^{1/2}(\Delta a_j+B_*u_j)$. Поскольку исходный остаток нулевой,
члены второго порядка в самой кривой не влияют на коэффициент при $t^2$
в квадрате остатка. Минимизация по $u_j$ даёт
\[
 u_j=-(B_*^\top C_jB_*)^{-1}B_*^\top C_j\Delta a_j.
\]
После подстановки коэффициент при $t^2$ равен
$\sum_j\pi_j(\Delta a_j)^\top S_j(\Delta a_j)$.
Вторая производная удваивает этот коэффициент.
\end{proof}

Положительная полуопределённость и точный ранг следуют из разложения
\begin{equation}\label{eq:schur-factor}
 S_j=C_j^{1/2}(I-P_{C_j^{1/2}B_*})C_j^{1/2}\succeq0,
 \qquad \rank S_j=\rank C_j-m.
\end{equation}
Здесь проектор евклидов, а равенство рангов использует
$B_*^\top C_jB_*\succ0$. Вырожденность $C_j$ сама по себе не запрещает
теорему. Но если $\rank C_j<m$, это условие невозможно; если
$\rank C_j=m$, окно даёт нулевую локальную кривизну профиля.
Нулевые окна $C_j=c_j=0$ можно удалить из суммы, сохранив остальные веса.
Подстановка псевдообратной в теорему без проверки гладкости профиля при
изменении ранга не обосновывает её расширение.

Пусть $V\in\R^{d\times(d-m)}$ --- ортонормированное дополнение $B_*$,
$\Delta=VK$. При векторизации по столбцам полная матрица формы есть
\begin{equation}\label{eq:tangent-matrix}
 H_* =\sum_j\pi_j(a_ja_j^\top)\otimes(V^\top S_jV)
 \in\R^{m(d-m)\times m(d-m)}.
\end{equation}
Она положительно определена тогда и только тогда, когда
\begin{equation}\label{eq:identification}
 \{K:S_j^{1/2}VK a_j=0\ \text{для всех }j\}=\{0\}.
\end{equation}
Это критерий \emph{невырожденной локальной идентифицируемости второго
порядка}, а не необходимое условие любой возможной идентифицируемости
по членам высшего порядка. Необходимая для $H_*\succ0$ счётная граница:
\[
 m(d-m)\le\sum_{j:a_j\ne0}(\rank C_j-m).
\]
Одной этой границы недостаточно: важны взаимные направления слагаемых.

\begin{remark}[Средней ковариации и возбуждения индексов недостаточно]
При $d=2,m=1$, $B_*=e_1$, $a_1=a_2=1$, $\pi_1=\pi_2=1/2$ возьмём
$C_1=u_1u_1^\top$, $C_2=u_2u_2^\top$, $u_1=(1,1)^\top$,
$u_2=(1,-1)^\top$. Тогда $\overline C=I_2$,
$\sum_j\pi_ja_j^2=1$, $B_*^\top C_jB_*=1$, но $S_1=S_2=0$.
Для любого близкого к $e_1$ направления $b$ можно подобрать
$g_j=(u_j^\top e_1)/(u_j^\top b)$ и получить нулевой остаток в каждом
окне. Полный профиль локально постоянен. Усреднение ковариаций и
демпфирование не создадут отсутствующую информацию.
\end{remark}

\section{Полная метрика: точный результат и границы}

В модели \eqref{eq:ideal} определим на горизонтальном пространстве
\[
 h_*(\Delta_1,\Delta_2)=
 \sum_j\pi_j(\Delta_1a_j)^\top S_j(\Delta_2a_j).
\]
Если $H_*\succ0$, это скалярное произведение и
$\Hess F_0(B_*)=2h_*$ как билинейные формы. Поэтому оператор
риманова гессиана в этой метрике равен $2I$ и имеет обусловленность $1$.
Это локальное утверждение в точном решении. Оно не означает один шаг
до решения, глобальную сходимость или устранение статистического шума.
При замене метрики в критической точке меняется оператор,
представляющий гессиан; его инерция и тип критической точки сохраняются.

Практическая подстановка текущих $B,g_j(B)$ вместо $B_*,a_j$ даёт
кандидат на предобусловливание. Чтобы это была гладкая риманова метрика,
нужна гладкость коэффициентов и матриц систем на рассматриваемой области,
а также положительная определённость формы. Добавление
$\epsilon\inner{\Delta_1}{\Delta_2}$ при $\epsilon>0$
обеспечивает положительную определённость, но само по себе не исправляет
негладкость псевдообратной на границе ранга.

В идеальной точке при $H_*\succ0$ собственные значения гессиана
относительно $H_*+\epsilon I$ равны
$2\lambda_i(H_*)/[\lambda_i(H_*)+\epsilon]$.
Если $H_*$ вырождена, добавление $\epsilon I$ в метрику не убирает
нулевые собственные значения гессиана. Малые значения кривизны могут
отражать слабую идентификацию: их обращение усиливает также шум.

\section{Для ridge: точный гессиан и две формы Гаусса--Ньютона}

Теперь $B$ --- произвольная текущая точка, остаток может быть ненулевым.
Предположим $A_j\succ0$, что обеспечивается $\tau>0$.
Для горизонтального $\Delta$ определим
\begin{equation}\label{eq:hessian-parts}
 z_j=\Delta g_j,\qquad q_j=\Delta^\top r_j,\qquad
 t_j=B^\top C_jz_j,\qquad k_j=t_j+q_j.
\end{equation}
Дифференцирование \eqref{eq:normal} даёт точную производную коэффициентов:
\begin{equation}\label{eq:coefficient-derivative}
 Dg_j(B)[\Delta]=-A_j^{-1}k_j.
\end{equation}

\begin{theorem}\label{th:ridge-hessian}
Для стандартной метрики Grassmann
$\langle\Delta_1,\Delta_2\rangle=\inner{\Delta_1}{\Delta_2}$
точная квадратичная форма гессиана профиля равна
\begin{equation}\label{eq:ridge-hessian}
 \boxed{\displaystyle
 \Hess F_\tau(B)[\Delta,\Delta]
 =2\sum_j\pi_j\left[
 z_j^\top C_jz_j+\tau\norm{z_j}^2-k_j^\top A_j^{-1}k_j
 \right].}
\end{equation}
\end{theorem}
\begin{proof}
Вдоль прямой $B+t\Delta$ теорема об огибающей и
\eqref{eq:coefficient-derivative} дают вторую производную
$2\sum_j\pi_j(z_j^\top C_jz_j-k_j^\top A_j^{-1}k_j)$.
Для полярной ретракции горизонтального направления
$B(t)=(B+t\Delta)(I+t^2\Delta^\top\Delta)^{-1/2}$
имеем $B''(0)=-B\Delta^\top\Delta$.
Добавка от ускорения кривой равна
\[
 \inner{G_E}{B''(0)}
 =-2\sum_j\pi_jr_j^\top B\Delta^\top\Delta g_j
 =2\tau\sum_j\pi_j\norm{\Delta g_j}^2,
\]
поскольку $B^\top r_j=-\tau g_j$.
Полярная ретракция имеет второй порядок для стандартной метрики
Grassmann, что завершает доказательство.
\end{proof}
В некритической точке гессиан другой переменной метрики может отличаться:
формула \eqref{eq:ridge-hessian} явно относится к стандартной метрике.

\paragraph{Форма после исключения линейных приращений коэффициентов.}
Положим
\begin{equation}\label{eq:ridge-schur}
 S_{j,\tau}=C_j-C_jB(B^\top C_jB+\tau I_m)^{-1}B^\top C_j\succeq0.
\end{equation}
Линеаризуем совместный расширенный остаток, а затем исключим приращения
локальных коэффициентов. Половина соответствующей формы Гаусса--Ньютона:
\begin{equation}\label{eq:schur-gn}
 h_{\rm SC}(\Delta,\Delta)
 =\sum_j\pi_j z_j^\top S_{j,\tau}z_j.
\end{equation}
Она всегда неотрицательна. Можно использовать
$h_{\rm SC}+\epsilon\norm{\Delta}_F^2$ как метрику с демпфированием.

\paragraph{Форма точного якобиана профилированного остатка.}
Введём
\[
 W_j=\begin{bmatrix}C_j^{1/2}B\\\sqrt\tau I_m\end{bmatrix},\quad
 b_j=\begin{bmatrix}y_j\\0\end{bmatrix},\quad
 e_j=b_j-W_jg_j,\quad W_j^\top e_j=0,
 \quad D_j=\begin{bmatrix}C_j^{1/2}\Delta\\0\end{bmatrix}.
\]
Точный якобиан переменного исключения линейных параметров имеет вид
\[
 De_j[\Delta]=-(I-P_{W_j})D_jg_j
              -W_jA_j^{-1}D_j^\top e_j.
\]
Два слагаемых ортогональны, $D_j^\top e_j=-q_j$. Поэтому половина
формы Гаусса--Ньютона \emph{точного профилированного остатка} равна
\begin{equation}\label{eq:varpro-gn}
 h_{\rm VP}(\Delta,\Delta)=
 h_{\rm SC}(\Delta,\Delta)+\sum_j\pi_jq_j^\top A_j^{-1}q_j.
\end{equation}
Терминология переменного исключения согласуется с \cite{varpro}.
Здесь формулы получены непосредственно для нашего ridge-профиля.

Точное отличие гессиана от \eqref{eq:schur-gn}:
\begin{equation}\label{eq:hessian-difference}
 \frac12\Hess F_\tau[\Delta,\Delta]-h_{\rm SC}(\Delta,\Delta)
 =\sum_j\pi_j\left[
 \tau\norm{z_j}^2-2t_j^\top A_j^{-1}q_j-q_j^\top A_j^{-1}q_j
 \right].
\end{equation}
При $\tau=0,r_j=0$ все три формы совпадают с теоремой
\ref{th:ideal}. При ridge или ненулевом остатке ни
$h_{\rm SC}$, ни $h_{\rm VP}$ нельзя без дополнительных условий называть
половиной точного гессиана. Последний может быть неопределённым, а формы
Гаусса--Ньютона остаются неотрицательными.

\section{Двусторонняя метрика и её действительная стоимость}

Пусть
\begin{equation}\label{eq:separable}
 M=\overline C+\eta I_d\succ0,\quad
 \overline C=\sum_j\pi_jC_j,\quad \eta>0,\qquad
 R_B=\sum_j\pi_jg_jg_j^\top+\epsilon I_m\succ0.
\end{equation}
Нормировка $M$ на среднее собственное значение допустима, но не меняет
его обусловленность; общий масштаб шага должен учитывать line search.
На горизонтальных векторах зададим
\begin{equation}\label{eq:separable-metric}
 h_B(\Delta_1,\Delta_2)
 =\tr[\Delta_1^\top S_M(B)\Delta_2R_B].
\end{equation}
Матрица $S_M$ имеет ядро $\spn(B)$ и положительна на его евклидовом
ортогональном дополнении. Поэтому \eqref{eq:separable-metric}
положительно определена на горизонтальном пространстве.
При $B\mapsto BQ$, $\Delta\mapsto\Delta Q$, $g_j\mapsto Q^\top g_j$
получаем $R_B\mapsto Q^\top R_BQ$, а значение формы сохраняется.
Это корректная метрика на Grassmann там, где $g_j(B)$ гладкие.

Через дополнение $V$ её градиент можно найти из
\[
 (V^\top S_MV)KR_B=V^\top G_E,\qquad \grad_hF=VK.
\]
Но явное дополнение и новая факторизация матрицы размера $d-m$
\emph{не требуются}.
\begin{proposition}\label{prop:inverse}
Для $V^\top V=I$, $V^\top B=0$ верно
\begin{equation}\label{eq:schur-inverse}
 (V^\top S_MV)^{-1}=V^\top M^{-1}V.
\end{equation}
Поэтому
\begin{equation}\label{eq:fast-step}
 \boxed{\displaystyle
 \grad_hF(B)=P M^{-1}P G_E R_B^{-1},\qquad
 \Delta=-\grad_hF(B).}
\end{equation}
\end{proposition}
\begin{proof}
В ортонормированном базисе $[B,V]$ матрица $V^\top S_MV$ является
дополнением Шура верхнего блока $B^\top MB$. Нижний диагональный блок
обратной матрицы равен обратному дополнению Шура, что даёт
\eqref{eq:schur-inverse}. Для произвольного горизонтального $Z$
подстановка \eqref{eq:fast-step} даёт
$h_B(\grad_hF,Z)=\inner{G_E}{Z}$.
\end{proof}

При фиксированном $M$ достаточно один раз выполнить Холецкого:
подготовка стоит $O(Jd^2+d^3)$. На шаге:
\[
 H=PG_E,\qquad MX=H,\qquad \grad_hF=PX R_B^{-1}.
\]
Один шаг сверх профиля и градиента стоит
\begin{equation}\label{eq:cost}
 O(Jm^2+d^2m+dm^2+m^3).
\end{equation}
Обращения в формулах реализуются решением систем.
$M$ можно фиксировать на всё внутреннее решение; $R_B$ пересчитывается
в принятой текущей точке. При $G_H\ne0$ направление является строгим
направлением спуска:
\[
 DF(B)[-\grad_hF]=-h_B(\grad_hF,\grad_hF)<0.
\]
Это утверждение о $F$, не об EDR.

Дальше применяется полярная или QR-ретракция с line search исходного
$F_\tau$. В частности, полярная ретракция:
$R_B(t\Delta)=(B+t\Delta)(I_m+t^2\Delta^\top\Delta)^{-1/2}$.
Не следует подменять возвращаемый \code{jac} у обычного SciPy L-BFGS
предобусловленным градиентом: \code{jac} обязан оставаться производной
переданного функционала. Нужны отдельный риманов алгоритм, корректное
предобусловливание внутри алгоритма либо согласованная замена координат.
Для L-BFGS на многообразии также нужны перенос касательных векторов
и согласованные секущие пары. Простейший проверочный алгоритм здесь ---
предобусловленный градиентный спуск с одинаковым Armijo во всех вариантах.

\section{Когда двусторонняя метрика имеет гарантию}

\begin{theorem}\label{th:separable}
В модели \eqref{eq:ideal} пусть $M\succ0$,
$\alpha M\preceq C_j\preceq\beta M$ для всех $j$,
$0<\alpha\le\beta$ и $G_*:=\sum_j\pi_ja_ja_j^\top\succ0$.
В метрике \eqref{eq:separable-metric} с $R_{B_*}=G_*$ собственные
значения гессиана в $B_*$ лежат в $[2\alpha,2\beta]$,
так что $\kappa\le\beta/\alpha$.
\end{theorem}
\begin{proof}
Для любой $C\succ0$
\[
 z^\top S_C(B_*)z=\min_u(z+B_*u)^\top C(z+B_*u).
\]
Из покомпонентного сравнения квадратичных форм и последующей
минимизации получаем
$\alpha S_M\preceq S_{C_j}\preceq\beta S_M$.
Применяем к $z=\Delta a_j$, суммируем с весами $\pi_j$ и используем
\eqref{eq:ideal-hessian}. Сумма с $S_M$ равна
$\tr(\Delta^\top S_M\Delta G_*)$.
\end{proof}

При $C_j=M$ все собственные значения равны $2$ независимо от
обусловленности $M$ и $G_*$. Если хотя бы одно $C_j$ вырождено,
условие $\alpha M\preceq C_j$ с $\alpha>0$ невозможно. Для приведённого
в разделе 1 сохранённого результата поэтому теорема не даёт гарантии.

Более подходящее условие при вырожденных окнах --- сравнение
\emph{совокупных} форм на касательном пространстве:
\begin{equation}\label{eq:aggregate-bound}
 \underline\lambda\,\mathcal P\preceq H_*\preceq
 \overline\lambda\,\mathcal P,\qquad
 \mathcal P=G_*\otimes(V^\top S_MV)\succ0.
\end{equation}
Тогда $\kappa\le\overline\lambda/\underline\lambda$.
Это следует из обобщённого спектра пары $(H_*,\mathcal P)$ и не требует
полноранговости каждого $C_j$. Но сама граница должна быть доказана
или измерена; из положительности $\overline C$ и $G_*$ она не следует.

Для текущего ridge-профиля вместо $H_*$ можно изучать формы
$h_{\rm SC}$ и $h_{\rm VP}$, а также точный гессиан
\eqref{eq:ridge-hessian}. Их спектры отвечают на разные вопросы.
В некритической точке сравнение стандартного гессиана с предобусловливателем
не является автоматически спектром гессиана переменной метрики.

\section{Статистический ridge: явное смещение и согласованный штраф}

\paragraph{Контрпример без ошибок моментов.}
Пусть одно окно, $d=2,m=1$, $B_*=e_1$,
\[
 C=\begin{pmatrix}2&1\\1&2\end{pmatrix},\qquad
 c=Ce_1=\begin{pmatrix}2\\1\end{pmatrix},\qquad\tau=1.
\]
Для единичного $b$ профилирование эквивалентно максимизации
$ (b^\top c)^2/[b^\top(C+I)b]$.
По обобщённому неравенству Коши--Буняковского оптимум достигается при
\[
 b_\tau\parallel(C+I)^{-1}c\parallel(5,1)^\top,
 \qquad \EDR(b_\tau,e_1)=1/\sqrt{26}\approx0.1961.
\]
При $\tau=0$ оптимум равен $\spn(e_1)$.
Таким образом, даже при точных моментах изменение метрики не может
заставить минимизатор обычного ridge совпасть с истиной: функционал
уже предпочитает другое подпространство.

\paragraph{Штраф, согласованный с локальной ковариацией.}
Другой статистический критерий получается при штрафе
\begin{equation}\label{eq:cov-penalty}
 \tau\norm{C_j^{1/2}Bg_j}^2
 =\tau (Bg_j)^\top C_j(Bg_j).
\end{equation}
Он инвариантен относительно любого обратимого изменения базиса и
штрафует энергию локального линейного предсказания. Параметр $\tau$
здесь безразмерный и не должен автоматически копироваться из масштаба
евклидова ridge.
\begin{proposition}\label{prop:cov-penalty}
Пусть $F^{\rm cov}_\tau$ --- профиль с \eqref{eq:cov-penalty},
а $E=\sum_j\pi_j\norm{y_j}^2$. При общих фиксированных $\tau\ge0$
и окнах для всех $B$ справедливо
\begin{equation}\label{eq:cov-affine}
 \boxed{\displaystyle F^{\rm cov}_\tau(B)
 =\frac{F_0(B)+\tau E}{1+\tau}.}
\end{equation}
Поэтому множество минимизирующих подпространств совпадает с $F_0$.
\end{proposition}
\begin{proof}
В каждом окне положим $D=C_j^{1/2}B$. Для произвольной матрицы $D$,
включая вырожденную, минимум
$\norm{y-Dg}^2+\tau\norm{Dg}^2$
достигается при $Dg=P_Dy/(1+\tau)$ и равен
$\norm{y}^2-\norm{P_Dy}^2/(1+\tau)$.
Нерегуляризованный минимум равен
$\norm{y}^2-\norm{P_Dy}^2$. Суммирование даёт результат.
\end{proof}

Этот штраф сохраняет исходное смещение нерегуляризованного функционала,
а не устраняет его. Он \emph{не стабилизирует идентификацию коэффициентов}:
локальная система $(1+\tau)B^\top C_jB$ имеет прежние ранг и
обусловленность. При разных $\tau_j$ меняются относительные веса окон,
и сохранение общего минимума в общем случае теряется. Если окна или
параметры зависят от $B$, формула остаётся поточечной, но вывод о тех же
минимизаторах из постоянного аффинного преобразования уже не следует.
Замена $C_j$ одной средней $M$ также не даёт \eqref{eq:cov-affine}.

\section{Связь ошибок моментов, ridge, точности решения и EDR}

Предыдущие результаты можно дополнить локальной \emph{детерминированной}
оценкой. Она не предполагает известную истину в рабочем алгоритме,
но её условия относятся к окрестности истины и должны обосновываться
отдельно. Это не вычислимый по одному fit сертификат и не
высоковероятностная статистическая теорема.

Пусть при фиксированных фактических $C_j$
\begin{equation}\label{eq:perturb}
 c_j=C_jB_*a_j+\xi_j,\qquad
 A_{*,j}=B_*^\top C_jB_*\succ0.
\end{equation}
В $\xi_j$ можно включить шум, ошибку локальной линейной модели и
эффект усечения моментов. Положим
\begin{equation}\label{eq:coefficient-bias}
 d_j=g_j(B_*)-a_j
 =(A_{*,j}+\tau I_m)^{-1}(B_*^\top\xi_j-\tau a_j),
 \qquad
 \norm{d_j}\le
 \frac{\norm{B_*^\top\xi_j}+\tau\norm{a_j}}
 {\lambda_{\min}(A_{*,j})+\tau}.
\end{equation}
Зафиксируем ортонормированное дополнение $V_*$ и полярную карту
\begin{equation}\label{eq:local-chart}
 B(K)=(B_*+V_*K)(I_m+K^\top K)^{-1/2},\qquad
 f(K)=F_\tau(B(K)).
\end{equation}
В нуле её производная равна $V_*K$, поэтому
\begin{align}\label{eq:score-bound}
 \nabla f(0)&=2\sum_j\pi_jV_*^\top(C_jB_*d_j-\xi_j)(a_j+d_j)^\top,
 \nonumber\\
 \norm{\nabla f(0)}_F&\le b_{\rm stat}:=
 2\sum_j\pi_j
 \left(\norm{V_*^\top C_jB_*}\norm{d_j}
       +\norm{V_*^\top\xi_j}\right)\norm{a_j+d_j}.
\end{align}
Матричные нормы без индекса в этой оценке операторные, векторные ---
евклидовы. Формула явно разделяет вклад ошибок моментов и вклад ridge.

\begin{theorem}[Локальная оценка EDR]\label{th:edr-bound}
Пусть на выпуклом шаре $\norm{K}_F\le\rho$ функция $f$ дважды
дифференцируема и $\nabla^2 f(K)\succeq\mu I$ с $\mu>0$.
Пусть вычисленная точка $\widehat K$ лежит в этом шаре и удовлетворяет
$\norm{\nabla f(\widehat K)}_F\le\zeta$. Тогда
\begin{equation}\label{eq:edr-bound}
 \boxed{\displaystyle
 \EDR(B(\widehat K),B_*)
 \le\frac{\norm{\widehat K}_F}{\sqrt m}
 \le\frac{b_{\rm stat}+\zeta}{\mu\sqrt m}.}
\end{equation}
\end{theorem}
\begin{proof}
Сильная монотонность градиента вдоль отрезка от $0$ до $\widehat K$ даёт
\[
 \mu\norm{\widehat K}_F^2\le
 \inner{\nabla f(\widehat K)-\nabla f(0)}{\widehat K}
 \le(\zeta+b_{\rm stat})\norm{\widehat K}_F.
\]
Отсюда вторая граница. Синусы главных углов в карте
\eqref{eq:local-chart} равны сингулярным значениям
$\widehat K(I+\widehat K^\top\widehat K)^{-1/2}$, так что
\eqref{eq:edr} даёт первую границу.
\end{proof}

Один достаточный способ получить $\mu$: пусть идеальный профиль с
теми же $C_j$ имеет матрицу \eqref{eq:tangent-matrix} с
$\lambda_{\min}(H_*)=\gamma>0$, его гессиан в карте меняется не более
чем на $L\rho$, а возмущение гессиана из-за $\xi_j,\tau$ ограничено
сверху числом $\delta_2$. Тогда
$\mu=2\gamma-L\rho-\delta_2>0$ является достаточной границей.
Здесь оценка $\delta_2$ является \emph{условием}, а не следствием одного
лишь малого остатка. Нужны также контроль локальных систем и гладкость.

В \eqref{eq:edr-bound} геометрия оптимизатора может уменьшить
$\zeta$ при данном бюджете. Она не убирает $b_{\rm stat}$ и не увеличивает
информацию $\gamma$ в данных. При малой $\gamma$ шум и смещение
усиливаются. Если остановка использует градиент в меняющейся метрике,
его малость следует переводить в ту же норму: в координатах карты при
$f$-метрике $\mathcal M(K)\preceq L_M I$ имеем
\[
 \norm{\nabla f(K)}\le\sqrt{L_M}\norm{\grad_{\mathcal M}f(K)}_{\mathcal M}.
\]
Поэтому одинаковый числовой порог в разных метрических нормах не означает
одинаковую точность оптимизации. Оценка требует попадания в шар и не
доказывает попадание в него из произвольной инициализации.

\section{Анизотропия ADP и зависимость метрики от точки}

Если ядро использует расстояние $\norm{Tx}^2=x^\top Ax$,
$A=T^\top T\succ0$, то естественная двойственная норма линейного
функционала $b^\top x$ равна
\begin{equation}\label{eq:dual}
 \sup_{x^\top Ax\le1}(b^\top x)^2=b^\top A^{-1}b.
\end{equation}
Поэтому $M=(T^\top T)^{-1}$ --- допустимый кандидат на метрику
коэффициентов. Если $T$ симметричен, $T^\top T=T^2$.
Такая форма ядра присутствует в \code{tex/data/multiindex.tex}.
Но двойственность норм не доказывает близость $M$ к кривизне
\eqref{eq:tangent-matrix}; это отдельная гипотеза. Регуляризация малых
собственных значений и масштабирование должны быть зафиксированы до
сравнения. Для $A\succ0$ имеем $\kappa(A^{-1})=\kappa(A)$.

\textbf{Важное исправление черновика.}
На обычном Grassmann разрешена гладкая метрика, зависящая от текущего
$B$. При вычислении \emph{градиента} в точке используются только значения
метрики в этой точке; её производные не нужны. Обычная QR- или полярная
ретракция остаётся ретракцией. Производные метрики появляются в
связности и точном римановом гессиане.

Другое дело --- ограничение $B^\top M(B)B=I_m$. Его дифференциал содержит
дополнительный член $B^\top DM(B)[\Delta]B$. Касательное пространство
и формулы \eqref{eq:generalized} для фиксированного $M$ тогда нельзя
применять без изменений. Фиксировать тензор на внутреннее решение удобно
для такого обобщённого ограничения, но это не обязательное требование
для переменной метрики на \emph{обычном} Grassmann.

\section{Как проверять гипотезу без смешения причин}

\paragraph{Отличие от уже проведённого исследования.}
В \code{experiments/full_moment_variants.py} вариант \code{metric}
умножает составные остатки и дизайн на \code{metric_roots}; он меняет
сам статистический критерий, а не только касательное направление.
В сохранённом отчёте
\code{verification/full-moment-variants-20261026/REPORT.md}
этот вариант снизил среднюю EDR с $0.32216$ до $0.30374$, но провалил
заранее заданный предел ухудшения одной пары: $+0.107565>0.08$.
Остальные три сравнения этого исследования также не прошли отбор,
проверочные серии не запускались. Здесь проверены отчёт и код
конструкции критерия; полный независимый пересчёт того исследования
заново не выполнялся. Эти результаты не являются проверкой метрики D
из \eqref{eq:separable-metric}, сохраняющей \eqref{eq:profile}.
Их нельзя выдавать ни за её успешную валидацию, ни за её опровержение.
Для нового сравнения нужны новые зафиксированные данные; старые gates
и отрицательные результаты сохраняются.

\paragraph{Сначала фиксированные окна.}
Для одного и того же $C_j,c_j,\pi_j,B_0,\tau$ сравнить одинаковый
алгоритм с одинаковым line search и разными метриками:
\begin{center}\small
\begin{tabular}{p{0.10\linewidth}p{0.76\linewidth}}
\toprule
Вариант & Метрика на обычном Grassmann\\\midrule
A & $\inner{\Delta_1}{\Delta_2}$\\
C & $\tr(\Delta_1^\top S_M\Delta_2)$\\
D & $\tr(\Delta_1^\top S_M\Delta_2R_B)$\\
E & $h_{\rm SC}+\epsilon\inner{\Delta_1}{\Delta_2}$;
    при необходимости отдельно $h_{\rm VP}+\epsilon\inner{\Delta_1}{\Delta_2}$\\
\bottomrule
\end{tabular}
\end{center}
Обобщённый Штифель B можно добавить как отдельную проверку представления,
с сохранённым штрафом \eqref{eq:invariant-ridge} и учётом
\eqref{eq:pullback}. Текущий QR-chart L-BFGS полезен как практический
контроль, но сравнение его с градиентным спуском D смешивает эффект
метрики и эффект алгоритма. Для причинной абляции A/C/D/E алгоритм
должен быть общим.

Для $d=25,m=5$ размер полного касательного оператора всего $100\times100$,
то есть $80\,000$ байт для одной плотной матрицы float64. Поэтому вначале
разумно построить его явно для диагностики, проверить спектр и сравнить
плотный шаг с матрично-свободным. Нельзя заранее объявлять такой solve
главным bottleneck. Само построение, локальные факторизации и повторные
оценки профиля также входят во время.

Приготовив $C_jB$, факторизации $A_j$ и $g_j$, применение оператора
$h_{\rm SC}$ к $\Delta$ можно вычислять через
\[
 P\sum_j\pi_j[S_{j,\tau}(\Delta g_j)]g_j^\top,
\]
не строя все $S_{j,\tau}$. Стоимость для плотных $C_j$:
$O(J(d^2+dm+m^2))$ на применение, дополнительно к подготовке.
Общая стоимость CG зависит от числа итераций и демпфирования.

\paragraph{Диагностика.}
Проверять значение \emph{того же} профиля, его стандартный горизонтальный
градиент, число вычислений, время подготовки и решения, ошибки локальных
систем, углы и ограничения. Записывать отдельно спектры $h_{\rm SC}$,
$h_{\rm VP}$ и, для малых задач, точного гессиана. Использовать $\tau=0$
только с явными условиями ранга; нулевые окна и границы ранга сохранять
в диагностике. Oracle-инициализация и малые повороты известной истины
годятся для синтетической проверки кривизны.

Уменьшение EDR на каждом шаге не является условием принятия шага и не
гарантируется. Например, при $C=I_2$, $c=e_1+\delta e_2$ и истине
$e_1$ минимум профиля направлен вдоль $c$; движение от истины к этому
минимуму уменьшает функционал и увеличивает истинную EDR.
На синтетике траекторию EDR можно сохранять как диагностику, но нельзя
использовать истину для выбора рабочего шага или настройки метрики.

\paragraph{Затем полный ADP.}
Окна, bandwidth и локальные моменты пересчитываются во внешнем цикле:
это уже другая, адаптивная проверка. Нужны одинаковые исходные данные,
начальные подпространства и случайные состояния, свежий фиксированный
парный отбор и нетронутая проверочная серия. Порог качества, бюджет,
правило выбора внешнего шага и демпфирование фиксируются заранее.
Следует отдельно учитывать EDR, долю восстановления, численные ошибки,
незавершённые fit, время процесса/fit/солвера и RSS. Настройки, изменяющие
штраф, исследуются отдельно от настроек метрики.
Изменение оптимизатора может изменить последующие окна и сам масштаб
следующего ridge; теоремы фиксированного профиля это не покрывают.

\section{Что доказано и что ещё предстоит проверить}

\begin{center}\small
\begin{tabular}{p{0.43\linewidth}p{0.46\linewidth}}
\toprule
Результат & Область действия\\\midrule
Неизменность допустимых подпространств при смене нормировки &
Фиксированная SPD-матрица; статистический критерий преобразован корректно\\
Идеальная полная форма = половина гессиана &
Нулевой остаток, $\tau=0$, $B_*^\top C_jB_*\succ0$\\
Единичная обусловленность полной метрики &
Дополнительно $H_*\succ0$, без демпфирования, только в точном решении\\
Граница для двусторонней метрики &
Теорема \ref{th:separable} или проверенное совокупное сравнение
\eqref{eq:aggregate-bound}\\
Точный ridge-гессиан и две GN-формы &
$A_j\succ0$, фиксированные окна и $\tau$\\
Оценка EDR через шум, ridge и точность решения &
Теорема \ref{th:edr-bound}; локальность и сильная выпуклость в карте
являются условиями\\
Выигрыш EDR/времени полного ADP от новой метрики &
В этой работе не измерялся и не доказан\\
\bottomrule
\end{tabular}
\end{center}

Практический порядок: сначала проверить $h_{\rm SC}$ и его спектр,
затем сопоставить с дешёвой D при общем алгоритме и неизменном профиле.
Обобщённый Штифель не нужен для применения D.
Анизотропный ridge представляет отдельную статистическую гипотезу.
Утверждать новизну конструкции без отдельного обзора нельзя:
предобусловливание Grassmann и обобщённого Штифеля уже исследовано
в \cite{boumal,shustin}; здесь изложен конкретный вывод для полных
локальных моментов.

\appendix
\section{Аудит исходного черновика}

\begin{enumerate}
\item Формулы обобщённого Штифеля, изменение физического ridge и
контрпример с анизотропным углом корректны при зафиксированных условиях.
Дополнена необходимая формула градиента сохранённого ridge.
\item Доказательство идеального гессиана корректно; добавлены точный ранг,
условие совокупной невырожденности и контрпример, когда усреднение
ковариаций скрывает отсутствие идентификации.
\item Подстановка оценённых $g_j$ не делает идеальную форму точным
гессианом действующего профиля. Выведены
\eqref{eq:ridge-hessian}--\eqref{eq:hessian-difference}.
\item Граница обусловленности разделённой метрики верна при указанных
предпосылках. Ограничение ранга $19$ из \code{N_loc=20} неверно;
применимость заменена проверкой реальных рангов и совокупного спектра.
\item Заявленная дополнительная кубическая стоимость на каждом шаге
завышена для фиксированного $M$: формула \eqref{eq:fast-step} использует
однократную факторизацию $M$.
\item Утверждение о необходимости производных переменной метрики для
самого градиентного шага неверно на обычном Grassmann. Оно заменено
различением метрики и зависящего от точки ограничения.
\item Добавлена локальная связь с обычной EDR. Её предпосылки не
объявляются автоматически выполненными в адаптивном ADP.
\end{enumerate}

\section{Воспроизводимая проверка формул}

Команда из корня репозитория:
\begin{quote}\ttfamily\small
OPENBLAS\_NUM\_THREADS=1 OMP\_NUM\_THREADS=1\\
python -m experiments.anisotropic\_geometry\_check
\end{quote}
Скрипт \code{experiments/anisotropic_geometry_check.py} использует
фиксированный seed $20261026$ и записывает
\code{verification/anisotropic-geometry-20261026/checks.json}.
Проверки включают конечные разности градиента и точного ridge-гессиана,
точный якобиан расширенного профилированного остатка, совпадение с
действующим прямым профилем, инвариантное продолжение ridge,
обращение двусторонней метрики, градиент/ретракцию обобщённого Штифеля,
его индуцированную метрику, идеальную кривизну вырожденных окон,
спектр при общей ковариации, границу $[2\alpha,2\beta]$,
контрпример ridge и тождество \eqref{eq:cov-affine}.

Это численная проверка алгебры на малых детерминированных задачах,
дополняющая приведённые доказательства. Она не заменяет независимый
математический аудит, не проверяет высоковероятностные условия
теоремы \ref{th:edr-bound} и не является экспериментом по восстановлению.

\begin{thebibliography}{9}
\bibitem{boumal}
N. Boumal, P.-A. Absil.
\emph{Low-rank matrix completion via preconditioned optimization on the
Grassmann manifold.} Linear Algebra and its Applications, 475, 200--239, 2015.
\url{https://optimization-online.org/wp-content/uploads/2012/09/3619.pdf}.

\bibitem{shustin}
B. Shustin, H. Avron.
\emph{Riemannian optimization with a preconditioning scheme on the
generalized Stiefel manifold.}
Journal of Computational and Applied Mathematics, 423, 114953, 2023.
\url{https://arxiv.org/abs/1902.01635}.

\bibitem{varpro}
D. M. O'Leary, B. W. Rust.
\emph{Variable Projection for Nonlinear Least Squares Problems.}
NIST, публикационная запись и первичный текст.
\url{https://www.nist.gov/publications/variable-projection-nonlinear-least-squares-problems}.
\end{thebibliography}
\end{document}
